\scanpage{274}
\mspage{P 30}

\deleted{``discrètes'' associées aux différents nombres premiers, devaient fournir la clef pour un renouvellement de vaste envergure de l'arithmétique. C'est dans cet esprit qu'il a dégagé, en 1949, les célèbres ``c o n j e c t u r e s \quad d e \quad W e i l''. Conjectures absolument époustouflantes, à vrai dire, qui faisaient entrevoir, pour ces nouvelles ``variétés'' (ou ``espaces'') de nature discrète, la possibilité de certains types de constructions et d'arguments (*)\unsure{,|\mbox{}} qui jusque là ne semblaient pensables que dans le cadre des seuls ``espaces'' considérés comme dignes de ce nom par les analystes - savoir, les espaces dits ``topologiques'' (où la notion de variation continue a cours).}

\deleted{On peut considérer que la géométrie nouvelle est,avant toute autre chose, une s y n t h è s e entre ces deux mondes, jusque là mitoyens et étroitement solidaires, mais pourtant séparés : le m o n d e \quad ``a r i t h m é t i q u e'', dans lequel vivent les (soi-disants) ``espaces'' sans principe de continuité, et le m o n d e \quad d e \quad l a \quad g r a n d e u r \quad c o n t i n u e, ou vivent les ``espaces'' au sens propre du terme, accessibles aux moyens de l'analyste et (pour cette raison même) acceptés par lui comme dignes de giter dans la cité mathématique. D a n s \quad l a \quad v i s i o n \quad n o u v e l l e, \quad c e s \quad d e u x \quad m o n d e s \quad j a d i s \quad s é p a r é s, \quad n ' e n \quad f o r m e n t \quad p l u s \quad q u ' u n \quad s e u l.}

\deleted{Le premier embryon de cette vision d'une ``g é o m é t r i e \quad a r i t h m é t i q u e'' (comme je propose d'appeler cette géométrie nouvelle) se trouve dans les conjectures de Weil. Dans le développement de certains de mes thèmes principaux (**) , ces conjectures sont restées ma principale source d'inspiration, tout au long des années entre 1958 et 1969. Dès avant moi, d'ailleurs, O s c a r \quad Z a r i s k i d'un côté, puis J e a n - P i e r r e \quad S e r r e de l'autre, avaient développé pour les espaces-sans-foi-ni-loi de la géométrie algébrique ``abstraite'' certaines méthodes ``topologiques'', inspirées de celles qui avaient cours précédemment pour les ``espaces bon teint'' de tout le monde (***).}

\rule{2cm}{0.4pt}

\deleted{(*) (A l'intention du lecteur mathématicien.) Il s'agit ici\unsure{.|\mbox{}} des ``constructions et arguments'' liés à la théorie cohomologique des variétés différentiables ou complexes, et notamment de ceux impliquant la formule des points fixes de Lefschetz, et la théorie de Hodge.}

\deleted{(**) Il s'agit des quatre thèmes ``médians'' (n\textsuperscript{o}s 5 à 8), savoir ceux des t o p o s de la c o h o m o l o g i e \quad é t a l e et \handwritten{$\ell$}-adique, des m o t i f s, et (dans une moindre mesure)\unsure{'|\mbox{}} celui des c r i s t a u x. J'ai dégagé ces thèmes tour à tour entre 1958 et 1966.}

\deleted{(***) (A l'intention du lecteur mathématicien.) La principale contribution de Zariski dans ce sens me paraît l'introduction de la ``topologie de Zariski'' (qui plus tard a été un outil essentiel pour Serre dans FAC), et son ``principe de connexité'' et ce qu'il a appelé sa ``théorie des fonctions holomorphes'' - devenus entre \deleted{\unsure{s|m}}mes mains la théorie des schémas formels, et les ``théorèmes de comparaison''}

% archive sheet number: 273
